%! TEX root = sebab_solusi_temu2.tex
% the main TeX file which is intended to compile, :VimtexReload after adjustment
% :h vimtex-tex-root
\documentclass[../main.tex]{subfiles}
% \includeonlyframes{current}
% \setbeameroption{hide notes} % Only slides
% \setbeameroption{show only notes} % Only notes
% \setbeameroption{show notes on second screen=right} % Both

\title{Identifikasi Risiko dan Prosedur Pencegahan Kecelakaan Kerja Konstruksi}
\author{Nama Pemateri}
\date{\today}

\begin{document}
\section{SMK3}%
\label{sec:SMK3}

% ---------------------------------------------------
% FRAME 15: Siklus Penerapan SMK3
% ---------------------------------------------------
\begin{frame}[mytitle=standard,c,mycolor=digiPH_uniblue!15]
	\frametitle{Siklus Penerapan SMK3}

	Berdasarkan Pasal 6 PP No. 50 Tahun 2012, SMK3 meliputi lima tahap:

	\begin{enumerate}
		\item \textbf{Penetapan Kebijakan K3}
		\item \textbf{Perencanaan K3}
		\item \textbf{Pelaksanaan Rencana K3}
		\item \textbf{Pemantauan dan Evaluasi Kinerja K3}
		\item \textbf{Peninjauan dan Peningkatan Kinerja SMK3}
	\end{enumerate}

	\begin{center}
		$\downarrow$\\
		\textit{Kembali ke kebijakan dan perencanaan}
	\end{center}

	\begin{exampleblock}{Artinya}
		SMK3 adalah siklus perbaikan berkelanjutan, bukan kegiatan satu kali.
	\end{exampleblock}
\end{frame}

% ---------------------------------------------------
% FRAME 16: Lima Tahap SMK3
% ---------------------------------------------------
\begin{frame}[mytitle=standard,t,mycolor=white]
	\frametitle{Lima Tahap SMK3: Dari Kebijakan hingga Perbaikan}

	\begin{table}
		\centering\small
		\begin{tabularx}{\textwidth}{|l|X|X|}
			\hline
			\textbf{Tahap} & \textbf{Pertanyaan kunci}           & \textbf{Contoh pada proyek jalan}        \\
			\hline
			1. Kebijakan   & Apa komitmen organisasi?            & Komitmen nol kecelakaan serius           \\
			\hline
			2. Perencanaan & Risiko apa yang harus dikendalikan? & Risiko alat berat, galian, lalu lintas   \\
			\hline
			3. Pelaksanaan & Bagaimana pengendaliannya?          & SOP, APD, rambu, barrier, briefing       \\
			\hline
			4. Pemantauan  & Apakah pengendalian efektif?        & Inspeksi, pengukuran, audit, investigasi \\
			\hline
			5. Peningkatan & Apa yang harus diperbaiki?          & Revisi metode kerja dan program K3       \\
			\hline
		\end{tabularx}
	\end{table}

	\begin{block}{Kunci: Plan $\rightarrow$ Do $\rightarrow$ Check $\rightarrow$ Improve}
		Kelima tahapan tersebut ditetapkan secara eksplisit dalam Pasal 6 PP No. 50 Tahun 2012.
	\end{block}
\end{frame}

% ---------------------------------------------------
% FRAME 17a: Studi Kasus dan Kesimpulan
% ---------------------------------------------------
\begin{frame}[mytitle=standard,t,mycolor=digiPH_creamy!15]
	\frametitle{Studi Kasus dan Kesimpulan}

	\begin{block}{Studi Kasus}
		Sebuah proyek peningkatan jalan melakukan pekerjaan: galian; pemadatan; penghamparan aspal; pengoperasian alat berat; dan pekerjaan pada lalu lintas aktif.

		\textbf{Pertanyaan analisis:}
		\begin{itemize}
			\item Apa saja bahaya utama? Bagaimana menentukan tingkat risikonya?
			\item Apa pengendalian yang harus diterapkan?
			\item Regulasi apa yang menjadi dasar?
			\item Bagaimana memastikan pengendalian tersebut efektif?
		\end{itemize}
	\end{block}

	\begin{exampleblock}{Kesimpulan}
		K3 konstruksi adalah kewajiban hukum sekaligus sistem manajemen. Proyek yang aman bukan proyek tanpa risiko, melainkan proyek yang mampu mengenali, mengendalikan, memantau, dan memperbaiki risikonya secara sistematis.
	\end{exampleblock}
\end{frame}

% ---------------------------------------------------
% FRAME 17b: Referensi dan Catatan
% ---------------------------------------------------
\begin{frame}[mytitle=standard,t,mycolor=white]
	\frametitle{Referensi dan Catatan}

	\begin{block}{Referensi Utama}
		\begin{itemize}
			\item UU No. 1 Tahun 1970 tentang Keselamatan Kerja.
			\item UU No. 2 Tahun 2017 tentang Jasa Konstruksi (diubah dgn UU 6/2023).
			\item PP No. 50 Tahun 2012 tentang Penerapan SMK3.
			\item UU No. 38/2004 tentang Jalan jo. UU No. 2/2022.
			\item PP No. 34 Tahun 2006 tentang Jalan.
			\item Permen PUPR No. 10 Tahun 2021 tentang Pedoman SMKK.
		\end{itemize}
	\end{block}

	\begin{alertblock}{Catatan Penting untuk Dosen}
		Untuk level sarjana akhir/pascasarjana, Slide 14–16 sebaiknya menjadi titik analisis utama: mahasiswa tidak hanya diminta menghafal regulasi, tetapi membedakan domain hukum K3 umum dan keselamatan konstruksi serta menerjemahkan 5 tahap SMK3 menjadi keputusan operasional proyek.
	\end{alertblock}
\end{frame}

\section{Resiko dan Pencegahannya}%
\label{sec:Resiko dan Pencegahannya}

% --- SLIDE 1 ---
\begin{frame}[mytitle=center, c, mycolor=digiPH_ocean, dark]
	\frametitle{Identifikasi Risiko dan Prosedur Pencegahan}
	\framesubtitle{Kecelakaan Kerja Konstruksi}

	\begin{center}
		\textbf{\Large Panduan Komprehensif Mewujudkan Lingkungan Kerja yang Aman (Zero Accident)}\\
		\vspace{0.8cm}
		\vspace{0.3cm}
	\end{center}

	\vspace{0.5cm}
	% \begin{block}{Saran Visual \& Prompt}
	% 	\scriptsize
	% 	\textbf{Visual:} Latar belakang proyek konstruksi modern saat matahari terbenam dengan siluet pekerja.\\
	% 	\textbf{Prompt Bing:} A cinematic wide shot of a modern construction site at sunset, featuring silhouettes of construction workers wearing hard hats and safety gear against a dramatic orange and purple sky, photorealistic, 4k.
	% \end{block}
\end{frame}

% --- SLIDE 2 ---
\begin{frame}[mytitle=standard, t]
	\frametitle{Pendahuluan}
	\framesubtitle{Mengapa K3 Konstruksi Vital?}

	\begin{itemize}
		\item \textbf{Definisi K3:} Keselamatan dan Kesehatan Kerja (K3) adalah pilar utama pelindung nyawa di sektor konstruksi.
		\item \textbf{Statistik Kecelakaan:} Sektor konstruksi menyumbang persentase tertinggi kecelakaan kerja global akibat sifat pekerjaannya yang dinamis dan berisiko tinggi.
		\item \textbf{Dampak Kecelakaan:} Mengakibatkan korban jiwa, cacat permanen, kerugian finansial, dan penundaan proyek.
	\end{itemize}

	\vfill
	% \begin{block}{Saran Visual \& Prompt}
	% 	\scriptsize
	% 	\textbf{Visual:} Grafik penurunan angka kecelakaan, didampingi foto pekerja yang sedang berdiskusi.\\
	% 	\textbf{Prompt Bing:} A group of diverse construction workers in bright high-visibility vests and yellow hard hats standing in a circle having a safety meeting on a sunny construction site, highly detailed, realistic.
	% \end{block}
\end{frame}

% --- SLIDE 3 ---
\begin{frame}[mytitle=imageplus, c, mycolor=digiPH_cream]
	\frametitle{Konsep Dasar}
	\framesubtitle{Bahaya vs. Risiko}

	\begin{itemize}
		\item \textbf{Definisi Bahaya (Hazard):} Segala sesuatu (sumber, situasi, tindakan) yang berpotensi menyebabkan cedera atau kerugian. \textit{Contoh: Kabel listrik terkelupas.}
		\item \textbf{Definisi Risiko (Risk):} Probabilitas atau kemungkinan bahaya tersebut benar-benar menyebabkan insiden, dikalikan dengan keparahan dampaknya. \textit{Contoh: Pekerja tersengat listrik.}
		\item \textbf{Tujuan Identifikasi:} Menemukan bahaya sebelum berubah menjadi kecelakaan.
	\end{itemize}

	\vfill
	% \begin{block}{Saran Visual \& Prompt}
	% 	\scriptsize
	% 	\textbf{Visual:} Ilustrasi perbandingan visual antara "Hazard" dan "Risk".\\
	% 	\textbf{Prompt Bing:} A clever split-screen safety illustration. Left side shows a heavy concrete block dangling from a crane (labeled Hazard). Right side shows a person standing directly underneath the dangling block (labeled Risk), flat vector art style, clean lines, educational.
	% \end{block}
\end{frame}

% --- SLIDE 4 ---
\begin{frame}[mytitle=standard, t, mycolor=digiPH_gray]
	\frametitle{Identifikasi Bahaya Utama}
	\framesubtitle{Fisik \& Mekanik}

	\begin{itemize}
		\item \textbf{Jatuh dari Ketinggian:} Risiko mematikan nomor satu (bekerja di atap, scaffolding tanpa guardrail).
		\item \textbf{Tertimpa Benda Jatuh (Struck-by):} Perkakas atau material yang jatuh dari lantai atas.
		\item \textbf{Bahaya Mekanik:} Terjepit mesin bertenaga besar, pergerakan ekskavator atau crane yang membahayakan pekerja di titik buta (\textit{blind spot}).
	\end{itemize}

	\vfill
	% \begin{block}{Saran Visual \& Prompt}
	% 	\scriptsize
	% 	\textbf{Visual:} Pekerja di ketinggian yang terpasang dengan aman pada struktur menggunakan full-body harness.\\
	% 	\textbf{Prompt Bing:} A close-up of a construction worker securely clipped to a high-rise steel beam using a heavy-duty full-body safety harness and twin lanyards, clear blue sky background, photorealistic, focus on the safety hooks.
	% \end{block}
\end{frame}

% --- SLIDE 5 ---
\begin{frame}[mytitle=standard, t]
	\frametitle{Identifikasi Bahaya}
	\framesubtitle{Kimia, Biologi, \& Ergonomi}

	\begin{itemize}
		\item \textbf{Bahaya Kimia:} Paparan debu silika dari pemotongan beton, uap cat beracun, dan bahan pelarut.
		\item \textbf{Bahaya Biologi:} Gigitan serangga di area proyek yang lembab, bakteri dari sanitasi yang buruk.
		\item \textbf{Ergonomi:} Mengangkat beban terlalu berat dengan postur yang salah, atau gerakan berulang yang menyebabkan cedera tulang belakang.
	\end{itemize}

	\vfill
	% \begin{block}{Saran Visual \& Prompt}
	% 	\scriptsize
	% 	\textbf{Visual:} Ilustrasi postur tubuh yang benar vs. salah saat mengangkat kotak berat.\\
	% 	\textbf{Prompt Bing:} 3D illustration showing the correct and incorrect ergonomic posture for lifting a heavy concrete block. One figure bending knees with a straight back (green checkmark), one figure bending at the waist (red X), industrial setting, clear and educational.
	% \end{block}
\end{frame}

% --- SLIDE 6 ---
\begin{frame}[mytitle=standard, c, mycolor=digiPH_cream]
	\frametitle{Metode Identifikasi Risiko}
	\framesubtitle{JSA \& HIRARC}

	\begin{itemize}
		\item \textbf{Job Safety Analysis (JSA):} Memecah satu tugas spesifik menjadi langkah-langkah kecil, mengidentifikasi bahaya di tiap langkah, dan menentukan solusinya.
		\item \textbf{HIRARC:} \textit{Hazard Identification, Risk Assessment, and Risk Control}. Pendekatan sistematis untuk memetakan seluruh aktivitas proyek dan memberi skor risiko (Rendah, Sedang, Tinggi).
		\item \textbf{Dokumentasi:} Semua temuan wajib dicatat dan disosialisasikan kepada pekerja sebelum pekerjaan dimulai.
	\end{itemize}

	\vfill
	% \begin{block}{Saran Visual \& Prompt}
	% 	\scriptsize
	% 	\textbf{Visual:} Seorang pengawas K3 (Safety Officer) memegang papan klip dan memeriksa area kerja.\\
	% 	\textbf{Prompt Bing:} A professional female construction safety inspector wearing a white hard hat and high-visibility jacket, holding a clipboard and pointing at a construction blueprint, checking the scaffolding structure in the background, realistic photography.
	% \end{block}
\end{frame}

% --- SLIDE 7 ---
\begin{frame}[mytitle=center, c, mycolor=digiPH_leaf, dark]
	\frametitle{Piramida Pengendalian Risiko}
	\framesubtitle{Hierarchy of Control}

	\begin{itemize}
		\item \textbf{Struktur Pengendalian:} Strategi 5 tingkat yang diurutkan dari yang paling efektif hingga paling tidak efektif.
		\item \textbf{Tingkatan:} Eliminasi (Paling Efektif) $\rightarrow$ Substitusi $\rightarrow$ Rekayasa $\rightarrow$ Administrasi $\rightarrow$ Alat Pelindung Diri (Paling Kurang Efektif).
		\item \textbf{Fokus K3:} Selalu upayakan pengendalian di tingkat teratas sebelum mengandalkan APD.
	\end{itemize}

	\vfill
	% \begin{block}{Saran Visual \& Prompt}
	% 	\scriptsize
	% 	\textbf{Visual:} Grafik piramida terbalik berwarna cerah yang menunjukkan 5 tingkat Hierarchy of Control.\\
	% 	\textbf{Prompt Bing:} An inverted triangle graphic showing the hierarchy of hazard controls in a construction theme. Five layers colored blue, green, yellow, orange, and red. Clean minimalist 3D rendering with small construction icons.
	% \end{block}
\end{frame}

% --- SLIDE 8 ---
\begin{frame}[mytitle=imageplus, t]
	\frametitle{Pengendalian Lini Pertama}
	\framesubtitle{Eliminasi \& Substitusi}

	\begin{itemize}
		\item \textbf{Eliminasi (Menghilangkan Bahaya):} Mengubah desain proyek sehingga pekerjaan berbahaya tidak perlu dilakukan. \textit{Contoh: Merakit komponen di tanah bawah, bukan di ketinggian.}
		\item \textbf{Substitusi (Mengganti):} Mengganti material atau proses berbahaya dengan yang lebih aman. \textit{Contoh: Mengganti cat berbahan dasar pelarut (solvent) dengan cat berbahan dasar air (water-based).}
		\item \textbf{Efektivitas:} Mencegah pekerja terekspos bahaya sama sekali.
	\end{itemize}

	\vfill
	% \begin{block}{Saran Visual \& Prompt}
	% 	\scriptsize
	% 	\textbf{Visual:} Drone pengawas bangunan digunakan alih-alih manusia yang memanjat struktur berbahaya.\\
	% 	\textbf{Prompt Bing:} A commercial drone flying close to inspect a damaged tall brick chimney on a construction site, blue sky, replacing the need for a human to climb, photorealistic, high tech industrial vibe.
	% \end{block}
\end{frame}

% --- SLIDE 9 ---
\begin{frame}[mytitle=standard, t, mycolor=digiPH_gray]
	\frametitle{Pengendalian Rekayasa}
	\framesubtitle{Engineering Controls}

	\begin{itemize}
		\item \textbf{Isolasi Bahaya:} Memisahkan bahaya dari pekerja secara fisik.
		\item \textbf{Contoh Fisik:} Memasang pagar pengaman (guardrails) di tepi gedung lantai atas, menutup lubang terbuka di lantai, dan memasang jaring penangkap (safety nets).
		\item \textbf{Sirkulasi \& Mesin:} Menggunakan sistem pembuangan debu pada alat potong keramik atau memasang sensor otomatis mati pada alat berat jika ada orang mendekat.
	\end{itemize}

	\vfill
	% \begin{block}{Saran Visual \& Prompt}
	% 	\scriptsize
	% 	\textbf{Visual:} Tepi bangunan bertingkat tinggi yang dilengkapi dengan pagar barikade jaring berwarna oranye yang kuat.\\
	% 	\textbf{Prompt Bing:} A high-rise building under construction showing strong bright orange safety netting and sturdy metal guardrails installed along the open edges of the concrete floors to prevent falls, realistic, daytime.
	% \end{block}
\end{frame}

% --- SLIDE 10 ---
\begin{frame}[mytitle=standard, c]
	\frametitle{Pengendalian Administratif}

	\begin{itemize}
		\item \textbf{Perubahan Cara Kerja:} Mengubah prosedur dan mengatur jadwal pekerja untuk meminimalkan durasi paparan risiko.
		\item \textbf{Rotasi Kerja:} Membatasi waktu pekerja berada di area bising agar tidak merusak pendengaran.
		\item \textbf{Standar Prosedur (SOP):} Memberlakukan sistem izin kerja (Permit to Work) untuk pekerjaan ruang terbatas (confined space).
		\item \textbf{Pelatihan:} Wajib mengikuti \textit{Toolbox Meeting} setiap pagi sebelum mulai bekerja.
	\end{itemize}

	\vfill
	% \begin{block}{Saran Visual \& Prompt}
	% 	\scriptsize
	% 	\textbf{Visual:} Rambu-rambu keselamatan peringatan bahaya di lokasi proyek.\\
	% 	\textbf{Prompt Bing:} A cluster of bright yellow and red construction safety warning signs attached to a wire mesh fence, reading various warnings like 'Danger', 'Hard Hat Area', and 'Authorized Personnel Only', realistic outdoor lighting.
	% \end{block}
\end{frame}

% --- SLIDE 11 ---
\begin{frame}[mytitle=imageplus, t, mycolor=digiPH_cream]
	\frametitle{Alat Pelindung Diri}
	\framesubtitle{APD / PPE}

	\begin{itemize}
		\item \textbf{Garis Pertahanan Terakhir:} Digunakan hanya ketika sisa bahaya tidak bisa dihilangkan oleh metode sebelumnya.
		\item \textbf{APD Wajib Standar:} Helm pelindung (hard hat), sepatu bot dengan pelindung baja (safety shoes), rompi reflektif, kacamata pelindung (safety glasses), dan sarung tangan.
		\item \textbf{APD Spesifik:} Masker respirator untuk area berdebu, earmuff untuk area bising, dan full-body harness untuk ketinggian.
		\item \textbf{Kelemahan:} Melindungi individu, bukan menghilangkan bahaya, dan bisa gagal jika rusak atau tidak dipakai dengan benar.
	\end{itemize}

	\vfill
	\begin{block}{Saran Visual \& Prompt}
		\scriptsize
		\textbf{Visual:} Potret jarak dekat pekerja proyek lengkap dengan APD standar.\\
		\textbf{Prompt Bing:} A highly detailed close-up portrait of a construction worker wearing a pristine yellow hard hat, protective safety goggles, bright orange high-visibility vest, and heavy duty work gloves, looking determined, industrial background.
	\end{block}
\end{frame}

% --- SLIDE 12 ---
\begin{frame}[mytitle=standard, t]
	\frametitle{Prosedur Tanggap Darurat}
	\framesubtitle{Emergency Response}

	\begin{itemize}
		\item \textbf{Kesiapsiagaan:} Rencana matang jika terjadi insiden (kebakaran, gempa, kecelakaan berat).
		\item \textbf{Fasilitas Darurat:} Lokasi kotak P3K (First Aid Kit), ketersediaan tandu, dan alat pemadam api ringan (APAR).
		\item \textbf{Jalur Evakuasi:} Titik kumpul (Muster Point) yang aman dan jalur evakuasi yang selalu bersih dari tumpukan material.
		\item \textbf{Komunikasi Darurat:} Daftar nomor kontak darurat (RS terdekat, pemadam kebakaran) yang terpampang jelas.
	\end{itemize}

	\vfill
	\begin{block}{Saran Visual \& Prompt}
		\scriptsize
		\textbf{Visual:} Papan informasi darurat dengan kotak P3K dan tandu di lokasi proyek.\\
		\textbf{Prompt Bing:} A designated emergency response station on a construction site featuring a bright green first aid kit mounted on a wall, a fire extinguisher, and a green sign pointing to the 'Muster Point', realistic detailed.
	\end{block}
\end{frame}

% --- SLIDE 13 ---
\begin{frame}[mytitle=standard, c, mycolor=digiPH_gray]
	\frametitle{Peran Pimpinan \& Partisipasi Pekerja}

	\begin{itemize}
		\item \textbf{Komitmen Manajemen:} Menyediakan anggaran K3, tidak mengorbankan keselamatan demi kecepatan proyek, dan memberi teladan disiplin.
		\item \textbf{Kewajiban Pekerja:} Mematuhi aturan K3, menggunakan APD, dan berani melaporkan kondisi tidak aman (Unsafe Condition) tanpa takut dihukum.
		\item \textbf{Budaya Keselamatan:} Mengubah K3 dari "kewajiban yang memberatkan" menjadi "budaya saling peduli".
	\end{itemize}

	\vfill
	\begin{block}{Saran Visual \& Prompt}
		\scriptsize
		\textbf{Visual:} Manajer proyek bersalaman dengan pekerja lapangan, sama-sama memakai perlengkapan keselamatan.\\
		\textbf{Prompt Bing:} A construction project manager in a white hard hat shaking hands with a general laborer in a yellow hard hat on site, expressing teamwork and mutual respect, bright daytime, realistic photography.
	\end{block}
\end{frame}

% --- SLIDE 14 ---
\begin{frame}[mytitle=standard, t]
	\frametitle{Studi Kasus}
	\framesubtitle{Keberhasilan vs Kegagalan}

	\begin{itemize}
		\item \textbf{Contoh Insiden (Kegagalan):} Runtuhnya scaffolding karena modifikasi ilegal oleh pekerja tanpa pengawasan, menyebabkan cedera parah.
		\item \textbf{Contoh Sukses (Keberhasilan):} Proyek X mencapai 1.000.000 jam kerja tanpa kecelakaan (Zero Accident) karena penerapan izin kerja ketat dan rotasi pengawasan harian.
		\item \textbf{Pelajaran:} Prosedur tertulis tidak berguna jika tidak ditegakkan dengan disiplin di lapangan.
	\end{itemize}

	\vfill
	\begin{block}{Saran Visual \& Prompt}
		\scriptsize
		\textbf{Visual:} Gambar area kerja yang sangat rapi dan bersih (Housekeeping yang baik).\\
		\textbf{Prompt Bing:} An exceptionally clean and well-organized indoor construction site, tools neatly stored on racks, no debris on the smooth swept floor, safety barriers correctly placed, bright industrial lighting, photorealistic.
	\end{block}
\end{frame}

% --- SLIDE 15 ---
\begin{frame}[mytitle=center, c, mycolor=digiPH_darkblue, dark]
	\frametitle{Kesimpulan}
	\framesubtitle{Sesi Tanya Jawab}

	\begin{itemize}
		\item \textbf{Ringkasan:} Identifikasi risiko yang proaktif dan penerapan Hierarchy of Control adalah kunci meminimalisir kecelakaan.
		\item \textbf{Slogan K3:} "Berangkat selamat, bekerja aman, pulang dengan selamat."
		\item \textbf{Call to Action:} Terapkan standar JSA hari ini di setiap area kerja Anda.
		\item \textbf{Sesi Q\&A:} Membuka ruang bagi audiens untuk bertanya dan berdiskusi.
	\end{itemize}

	\vfill
	% \begin{block}{Saran Visual \& Prompt}
	% 	\scriptsize
	% 	\textbf{Visual:} Helm proyek dengan tanda tanya atau ikon mikroskopis yang menyimbolkan analisis mendalam.\\
	% 	\textbf{Prompt Bing:} A striking conceptual image of a yellow hard hat resting on a wooden table, with a glowing, semi-transparent 3D question mark hovering slightly above it, cinematic lighting, shallow depth of field.
	% \end{block}
\end{frame}

\end{document}

